\section{The exact additive quotient projector}
\label{sec:tpc39-projector}

The multiplicative CRT coordinates have already done their job once
\eqref{eq:tpc39-interface-recombination} is reached.  The surviving
problem is finite additive recovery on the alias index.

\subsection{The minimal no--wrap quotient}

Retain the notation in \eqref{eq:tpc39-interface-alias-ledger} and put
\begin{equation}
 q=L+1,
 \qquad
 \omega=\ee{1/q}.
 \label{eq:tpc39-projector-q}
\end{equation}
The integer \(q\) need not be prime.

\begin{lemma}[Minimal quotient size]
\label{lem:tpc39-qM-window}
One has
\begin{equation}
 B<qM\le B+M.
 \label{eq:tpc39-qM-window}
\end{equation}
Moreover, the residue class \(0\pmod q\) meets
\(\{-L,\ldots,L\}\) only at \(0\).
\end{lemma}

\begin{proof}
Write \(B=LM+r\), where \(0\le r<M\).  Then
\((L+1)M=B+(M-r)\), proving \eqref{eq:tpc39-qM-window}.  Since
\(q=L+1\), every nonzero \(k\) in the alias interval satisfies
\(0<|k|<q\).
\end{proof}

\begin{theorem}[Additive quotient identity]
\label{thm:tpc39-additive-quotient-identity}
For every integer \(F\) with \(|F|\le B\),
\begin{equation}
 \boxed{
 \one_{M\mid F}\frac1q\sum_{\nu=0}^{q-1}
 \ee{\nu F/(qM)}
 =\one_{qM\mid F}
 =\one_{F=0}.}
 \label{eq:tpc39-additive-quotient-identity}
\end{equation}
\end{theorem}

\begin{proof}
If \(M\nmid F\), the first member is zero.  If \(F=kM\), additive
orthogonality gives
\[
 \frac1q\sum_{\nu=0}^{q-1}\omega^{\nu k}
 =\one_{q\mid k}.
\]
Thus the first member is \(\one_{qM\mid F}\).  By
\cref{lem:tpc39-qM-window}, the only such determinant in the window is
\(F=0\).
\end{proof}

The first factor in \eqref{eq:tpc39-additive-quotient-identity} is
essential in the physical construction.  The additive average by itself
is generally nonzero when \(M\nmid F\), so it is not a divisibility
indicator on the full determinant window.  The product with
\(\one_{M\mid F}\) is the indicator \(\one_{qM\mid F}\).  The available
completed objects are the separate full and proper face outputs; their
nonjoint terms disappear only after using
\eqref{eq:tpc39-interface-recombination} with the same twist on both
fields.

\subsection{Twisted face outputs}

For \(0\le\nu<q\), define
\begin{align}
 \mathscr Y_\nu^C
 &:=%
 \sum_{|F|\le B}
 \mathscr Z_F\!\left[
  \ee{\nu F/(qM)}C
 \right],
 \label{eq:tpc39-Ynu-C}\\
 \mathscr Y_\nu^P
 &:=%
 \sum_{|F|\le B}
 \mathscr Z_F\!\left[
  \ee{\nu F/(qM)}P
 \right],
 \label{eq:tpc39-Ynu-P}\\
 \mathscr Y_\nu
 &:=\mathscr Y_\nu^C+\mathscr Y_\nu^P.
 \label{eq:tpc39-Ynu}
\end{align}

\begin{proposition}[Recombined twist formula]
\label{prop:tpc39-recombined-twist}
For every \(\nu\),
\begin{equation}
 \boxed{
 \mathscr Y_\nu
 =\sum_{k=-L}^{L}\omega^{\nu k}\mathscr Z_k.}
 \label{eq:tpc39-recombined-twist}
\end{equation}
In particular, all nonjoint face atoms cancel before the alias bank is
formed.
\end{proposition}

\begin{proof}
Linearity and \eqref{eq:tpc39-interface-recombination} give
\[
 \mathscr Y_\nu
 =\sum_{|F|\le B}
 \ee{\nu F/(qM)}\one_{M\mid F}\mathscr R_F.
\]
Write \(F=kM\) and use \eqref{eq:tpc39-interface-Zk}.
\end{proof}

\begin{theorem}[Exact Hilbert equality recovery]
\label{thm:tpc39-exact-Hilbert-recovery}
For arbitrary columns \(\mathscr Z_k\in\cH_{\rm phys}\),
\begin{equation}
 \boxed{
 \mathscr Z_0
 =\frac1q\sum_{\nu=0}^{q-1}\mathscr Y_\nu.}
 \label{eq:tpc39-exact-Hilbert-recovery}
\end{equation}
Consequently,
\begin{equation}
 \boxed{
 \|\mathscr Z_0\|^2
 \le\frac1q\sum_{\nu=0}^{q-1}\|\mathscr Y_\nu\|^2.}
 \label{eq:tpc39-condition-one-recovery}
\end{equation}
\end{theorem}

\begin{proof}
Average \eqref{eq:tpc39-recombined-twist} over \(\nu\).  By
\cref{lem:tpc39-qM-window}, only \(k=0\) survives.  Jensen's inequality
gives \eqref{eq:tpc39-condition-one-recovery}.
\end{proof}

There is no factor \(q\) in
\eqref{eq:tpc39-condition-one-recovery}.  The recovery functional is the
orthogonal average of a regular polygon bank.  The later minimality theorem
shows that this sample count and normalized recovery amplification are
simultaneously optimal.  This is not a condition-number assertion for the
full tomography map, which has a nontrivial kernel.

\begin{remark}[What has and has not been recovered]
\label{rem:tpc39-algebra-versus-analysis}
The equality column is recovered exactly from the \emph{values} of all
twisted completed outputs.  No bound for those outputs has yet been
proved.  Equation \eqref{eq:tpc39-condition-one-recovery} converts a
twist--stable energy theorem into the desired physical estimate; it is
not itself such an energy theorem.
\end{remark}
